\section{Empirical Evaluation}
\label{section:evaluation}

\subsection{Performance Metrics}
We plan to evaluate the algorithms based on the following four metrics:
\begin{itemize}
    \item{Final validation performance:} We record the best validation loss obtained by the model trained with the hyperparameters discovered by the agent, using the task-appropriate metric (e.g.\ RMSE for regression, log-loss for classification).
    \item{Search efficiency:} This metric evaluates how quickly the agent converges to strong configurations. We quantify it by the number of function evaluations (the evaluation budget), reporting performance as a function of the budget so that cold-start and asymptotic behavior can be read off a single anytime curve.
    \item{Simple regret:} Derived from the multi-armed bandit literature, simple regret measures the loss gap between the best configuration found by any compared method (the oracle) and the agent's recommended configuration. Because the benchmark spans widely varying loss scales, we report \emph{normalized} simple regret to make the metric comparable across tasks.
\end{itemize}

\subsection{Baseline Methods}
We compare against the following standard HPO baselines:
\begin{itemize}
    \item{Random Search} \citep{bergstra2012random}: A fundamental baseline that randomly samples hyperparameter configurations. It serves as a necessary lower bound to verify that the proposed RL agent explores high-dimensional continuous search spaces more effectively than random selection.
    \item{BO} \citep{snoek2012practical}: A standard state-of-the-art method for HPO. We will use a Gaussian Process-based BO with Expected Improvement (EI) to systematically benchmark against our RL agent's sequential decision-making capabilities.
    % \item{Hyperband} \citep{li2017hyperband}: An efficient, bandit-based approach that allocates training resources dynamically. It uses successive halving to aggressively early-stop underperforming deep learning models, serving as a strong baseline for search efficiency and resource allocation.
    \item{Hyp-rl} \citep{jomaa2019hyp}: An RL-based HPO method that learns a sequential decision policy via DQN over a discretized hyperparameter space, using only scalar reward feedback without learning curve information.
\end{itemize}

In addition, we evaluate two variants of our proposed method:
\begin{itemize}
    \item \textbf{CrossDataset-HyperRL}: A Hyp-RL agent pre-trained across multiple 
    source datasets and evaluated on held-out target datasets, isolating the 
    contribution of cross-dataset policy transfer (idea ii).
    \item \textbf{CrossDataset-LC-DQN}: Our full method, which combines cross-dataset 
    pre-training with the derivative-informed state representation (idea i), 
    augmenting the agent's input with first- and second-order derivatives of the 
    observed learning curves.
\end{itemize}

\subsection{Benchmark Tasks or Datasets}
We evaluate the algorithms in various benchmark domains widely used in the literature, specifically focusing on continuous prediction tasks:
\begin{itemize}
    \item \textbf{LCBench} \citep{zimmer2021autopytorch}: A learning curve benchmark built on top of OpenML datasets, providing pre-computed training trajectories of funnel-shaped MLPs across 35 tabular classification tasks. Each of its 2{,}000 hyperparameter configurations per dataset is logged with per-epoch validation losses over 50 epochs, exposing the full learning curve required by our derivative-informed state representation. The shared 7-dimensional search space (batch size, learning rate, momentum, weight decay, number of layers, max units per layer, dropout) and aligned training protocol across datasets make LCBench particularly well-suited for evaluating cross-dataset policy transfer.
    % \item \textbf{NAS-Bench-360} \citep{tu2021nasbench360}: While traditional NAS benchmarks focus on classification, NAS-Bench-360 extends to diverse domains including 1D sequence and continuous prediction. We plan to utilize this benchmark in future work to evaluate the RL agent's capability in navigating discrete structural hyperparameter spaces (e.g., layer types, channel sizes) for regression models.
    \item \textbf{Kaggle Playground Series} \citep{kaggle_playground}: To evaluate the generalization and cross-dataset transfer capabilities of the reinforcement learning policy, we utilize 15 tabular datasets sourced from the Kaggle Playground Series competitions. This benchmark suite encompasses variations in feature dimensions, instance counts, and mathematical task types (spanning continuous regression and discrete classification). The datasets, identified by their competition season and episode (e.g., \texttt{playground-series-s3e1}), consist of synthetic data generated via deep learning models trained on real-world reference data. This generation mechanism ensures the preservation of underlying statistical distributions and feature correlations, providing a standardized, heterogeneous environment for evaluating policy transferability.
\end{itemize}

% \bibliographystyle{plainnat}
% \bibliography{refs}